\section{Discussion}\label{sec:discussion}

\subsection{Classification of results}

The ten problems in the First Proof benchmark span a remarkably broad
swath of contemporary mathematics.  We organize our results into
three tiers:

\medskip
\noindent\textbf{Tier 1: Complete solutions (8 problems).}
Problems Q1, Q2, Q3, Q5, Q7, Q8, Q9, and Q10 are fully resolved.
These problems admit clean, self-contained proofs that, while requiring
substantial domain expertise, do not encounter fundamental obstructions.
The proof techniques range from:
\begin{itemize}[nosep]
\item \emph{External deep theorems} (Q1: Hairer's singularity result for $\Phi^4_3$;
  Q2: Bernstein--Zelevinsky derivative theory;
  Q7: Farrell--Jones conjecture + Selberg's lemma + $L$-theory transfer);
\item \emph{Explicit counterexamples} (Q3: direct computation at $q=1$
  with specific partitions);
\item \emph{Systematic constructions} (Q5: subgroup-order induction on
  geometric fixed points; Q8: cotangent-graph mollification;
  Q9: mode-separation and 5$\times$5 minors; Q10: matrix-free PCG
  with Kronecker preconditioner).
\end{itemize}

\medskip
\noindent\textbf{Tier 2: Substantial partial results (2 problems).}
Problems Q4 and Q6 have extensive partial solutions that identify the
precise remaining obstacle:
\begin{itemize}[nosep]
\item \textbf{Q4} ($\boxplus_n$--$\Phi_n$ inequality):
  the inequality is proved for $n=2$ (exact computation), $n=3$
  (Cauchy--Schwarz and Jensen), and for all degenerate/monotone
  configurations at arbitrary $n$.
  The general $n\ge 4$ case is reduced to a conjecture about the
  monotonicity of $1/\Phi_n$ under the heat flow $\partial_t p = \tfrac12 p''$,
  which would follow from a de~Bruijn-type identity for the $\Phi_n$
  functional.  This appears tractable but requires new analytic estimates.

\item \textbf{Q6} ($\varepsilon$-light sets):
  we establish an effective-resistance normalization (Lemma~4),
  a Foster-theorem-based bound on dangerous edges (Lemma~7),
  a Caro--Wei independent set of size $\Omega(\varepsilon n)$ in the
  dangerous-edge graph (Proposition~8), and complete proofs for seven
  special graph families.
  The general case is reduced to a single ``vertex BSS lemma''---a
  spectral partition result for rank-one matrices with bounded norms---which
  is a natural extension of the Marcus--Spielman--Srivastava
  theorem~\cite{mss2015} to vertex-induced subsets.
\end{itemize}

\medskip
\noindent\textbf{Note on Q7 (lattices with 2-torsion).}
This problem, touching upon deep questions in geometric topology,
required the most substantial novel argument.  The solution has two
logically independent components:
\begin{itemize}[nosep]
\item \textbf{NO} for lattices with odd-order torsion (Fowler's obstruction).
\item \textbf{Unconditional YES} for all lattices with only 2-torsion and
  $d=\dim(G/K)\ge 5$, covering \emph{all residues of $d$ modulo~$4$}
  and \emph{all} connected semisimple $G$ (both equal-rank and non-equal-rank).
  The key innovation is a \emph{transfer vanishing lemma} (Lemma~7.3):
  using Selberg's torsion-free subgroup $\Gamma'\le\Gamma$
  and the induction--restriction formula $\mathrm{ind}\circ\mathrm{res}
  = [\Gamma:\Gamma']\cdot\mathrm{id}$ in $L$-theory,
  one shows $[\Gamma:\Gamma']\cdot s(B\Gamma)=0$ in
  $\mathbb S_d(B\Gamma)$, hence $s(B\Gamma)\otimes\mathbb Q=0$.
  This bypasses the failure of the $B\Gamma$-assembly to be a rational
  isomorphism in dimensions $d\equiv 0\pmod{4}$.
\end{itemize}

\subsection{Methodological observations}

Several patterns emerged across the ten problems:

\begin{enumerate}[nosep]
\item \textbf{Literature leverage is decisive.}
  Problems Q1 and Q2 are solvable primarily because deep theorems
  (Hairer 2022; Bernstein--Zelevinsky theory) provide the heavy lifting.
  The proof-writing task is then to correctly \emph{assemble} these results
  into a coherent argument.  This is precisely the ``final stage'' of
  mathematical research that the First Proof benchmark targets.

\item \textbf{Counterexamples require search, not construction.}
  Q3 is answered by a counterexample.  Finding it required systematic
  exploration of the parameter space $(n, \lambda, t)$ at $q=1$,
  guided by structural understanding of what makes the stationary
  distribution ``non-Markov-compatible.''

\item \textbf{Partial results sharpen the frontier.}
  For Q4 and Q6, the partial solutions are arguably more valuable
  than a simple ``open'' verdict.  Each identifies a precise, well-formulated
  conjecture or lemma that, if proved, would close the problem.
  This kind of ``gap identification'' is itself a research contribution.
  Q7, initially appearing similarly intractable, was ultimately resolved
  in full by the transfer vanishing lemma---illustrating how a single
  conceptual insight can close a problem that resisted direct attack.

\item \textbf{Proof architecture matters.}
  Q8 (polyhedral Lagrangian smoothing) required multiple iterations to
  reach a logically closed proof.  The final 5-step architecture
  (Euler obstruction $\to$ cotangent-graph charts $\to$ mollification
  $\to$ global assembly $\to$ Hamiltonian isotopy) emerged only after
  several false starts involving tropical geometry and vertex resolution
  that, while conceptually appealing, introduced unnecessary complexity.
\end{enumerate}

\subsection{Comparison with the First Proof benchmark design}

The benchmark paper~\cite{firstproof} notes that each question
``has been solved by the author(s) of the question with a proof that is
roughly five pages or less.''  Our solutions range from approximately
1 page (Q1, Q3) to 8--10 pages (Q4, Q6, Q8), with the longer solutions
reflecting either partial results with extensive case analysis or
multi-step constructions with all intermediate lemmas proved in full.

The benchmark was designed to test whether AI systems can solve problems
``on their own, without an expert in the loop.''
Our approach explicitly includes human review in the loop, which we view
as a more realistic model of how AI will be used in mathematical research:
not as an autonomous solver, but as a collaborator whose output is
verified and refined by human judgment.

A distinctive feature of our work is the inclusion of Lean~4 formal
proof skeletons for all eight complete solutions.
While the formalizations necessarily axiomatize deep external theorems
(Hairer's singularity result, Bernstein--Zelevinsky derivative theory,
Blumberg--Hill equivariant framework, the Farrell--Jones conjecture,
etc.), they verify the logical
deduction chains that constitute our original contributions.
For Q7 in particular, the Lean skeleton formalizes the entire
transfer vanishing argument (Lemma~7.3) and the deduction chain from
Selberg's lemma through the induction--restriction formula to the
surgery realization, with no \texttt{sorry}.
This represents a step toward the vision articulated
by the Project Omega framework~\cite{omega-theory}: that mathematical
proofs should be \emph{multi-layered}, passing through human understanding,
peer review, and machine verification.
We note that the counterexample for Q3 is the most fully formalized,
with exact rational arithmetic in Lean~4 verifying all numerical evaluations.

\subsection{Open problems}

Two well-defined open problems emerge from our work:

\begin{enumerate}[nosep]
\item \textbf{Q4, general $n\ge 4$:}
  Prove or disprove that $1/\Phi_n(p\boxplus_n q)\ge 1/\Phi_n(p)+1/\Phi_n(q)$
  for all monic real-rooted polynomials of degree $n\ge 4$.
  The conjectured approach via heat-flow monotonicity of $\Phi_n$
  is outlined in Section~\ref{sec:q4}.

\item \textbf{Q6, vertex BSS lemma:}
  Given a connected graph $G$, $\varepsilon\in(0,1]$, and a vertex set $I$
  of size $\Omega(\varepsilon n)$ with all internal edge resistances $\le\varepsilon$,
  does there exist $S\subseteq I$ with $|S|=\Omega(|I|)$ and
  $M(S)\preceq \varepsilon I$?
\end{enumerate}

\noindent
Note that Q7 (pure 2-torsion lattices), previously listed as open for
$d\equiv 0\pmod{4}$, has been completely resolved by the
transfer vanishing lemma (Lemma~7.3 in Section~\ref{sec:q7}).
